\section*{Résumé}

Ce mémoire présente la conception d'une plateforme de formation autonome alimentée par
l'intelligence artificielle, développée au sein de l'entreprise Audit Énergie. Initialement
positionnée sur le courtage en énergie, l'entreprise s'est récemment structurée comme un
organisme de formation externe préparant des apprenants à des titres professionnels. Pour
une petite structure, ce développement se heurte à une contrainte forte : le coût d'un
formateur en direct rend le modèle difficilement viable, tandis que les cours simplement
préenregistrés appauvrissent l'accompagnement pédagogique.

Le travail propose de répondre à ce problème par une plateforme capable d'automatiser la
production, la diffusion et l'accompagnement des cours. Deux contributions principales sont
étudiées. La première est une \emph{pipeline} de génération de contenus pédagogiques, qui
transforme un référentiel officiel en cours structurés, oralisables et calibrés en durée. Plutôt
qu'une génération directe par un prompt unique, cette pipeline s'appuie sur un plan
intermédiaire explicite, une génération section par section, des contrôles qualité ciblés et la
conservation d'artefacts intermédiaires garantissant la traçabilité. La seconde contribution est
un assistant pédagogique de type RAG (\emph{Retrieval-Augmented Generation}), fondé sur
Azure AI Search et Azure OpenAI, qui répond aux questions des élèves à partir du contenu réel
de la formation plutôt que d'une connaissance générale non contrôlée.

Le mémoire situe ces choix dans l'état de l'art (prompt engineering, fine-tuning, RAG,
génération structurée, reviews automatiques) et discute leurs limites. Il insiste sur une
exigence centrale : encadrer l'IA générative pour la rendre fiable, auditable et mesurable dans
un usage réel de production. Sont enfin présentées les conditions d'application de l'approche,
ses limites et les indicateurs d'évaluation associés, qu'il s'agisse de l'efficacité opérationnelle,
de la qualité pédagogique ou de la pertinence des réponses de l'assistant.

\vspace{0.4cm}
\noindent\textbf{Mots-clés :} intelligence artificielle générative, formation à distance,
pipeline de génération de contenus, synthèse vocale, RAG, Azure AI Search, traçabilité,
automatisation pédagogique.
